\begin{helper}[The differentiated equation (3.2)]
Let $n\ge2$, $p>1$, $\varepsilon>0$ and $0<\gamma_0\le\gamma_1$.  Let
$\Omega\subseteq\R^n$ be open, let $A\colon\R^n\to\R^n$ satisfy
(Q1) of the thesis--(Q2) of the thesis on $\Omega$ with these parameters
(\noderef{EllipticStructureQ12}), and assume the cited regularity result (C1)
for $\varepsilon>0$ (\noderef{CitedRegularityEpsPos}).  Let $u$ be a weak
solution of $\operatorname{div}A(\nabla u)=0$ in $\Omega$
(\noderef{IsWeakSolutionOn}) with weak gradient $g=\nabla u$, and let
$Dg^k$ be the weak gradient of the $k$-th component $g^k=D_ku$, so that
$u\in W^{2,2}_{\mathrm{loc}}(\Omega)$ with second derivatives
$D_jD_ku=(Dg^k)^j$ (\noderef{MemW1pLocOn}).  Then for every
$k\in\{1,\dots,n\}$ and every $\varphi\in W^{1,2}(\Omega)$ with weak gradient
$g_\varphi$ (\noderef{MemW1pOn} at the exponent $2$) and with compact support
contained in $\Omega$,
\[
\int_\Omega\Bigl\{\sum_{i,j=1}^nD_{h_j}A^i(\nabla u)\,D_jD_ku\,D_i\varphi\Bigr\}dx=0 .
\]

\textbf{On the test class.}  The thesis prints (3.2) of the thesis ``for all
$\varphi\in W^{1,p}(\Omega)$ with $\operatorname{supp}\varphi$ compact,
$\operatorname{supp}\varphi\subset\Omega$'' (p.~16), copying the test class of
(1.2) of the thesis.  The class above is $W^{1,2}$ with compact support, which for
$p>2$ is strictly \emph{larger}, so the statement above is stronger than the
printed one.  That strengthening is forced by the paper itself, in the following
precise sense.

The thesis prints the $W^{1,p}$ class immediately after assuming
$u\in C^2$ (p.~16), and the Note for the formalization deletes that assumption,
leaving $u\in W^{2,2}_{\mathrm{loc}}\cap C^1$.  Under the weaker regularity the
printed class no longer contains the test functions the thesis itself
substitutes into (3.2) of the thesis twenty lines later: with $w=\varepsilon+|\nabla
u|^2$ and $\zeta\in C_0^\infty(\Omega)$, the gradients of
$\varphi=D_ku\,w^{t/2}\zeta^2$ (p.~16, giving (3.3) of the thesis) and of
$\varphi^{\pm}=\pm(D_ku-\ell)^{\pm}\zeta^2$ (p.~16, giving (3.4) of the thesis)
contain $\nabla(D_ku)$, which the Note's regularity places in
$L^2_{\mathrm{loc}}(\Omega)$ and in nothing better; so for $p>2$ those two
functions lie in $W^{1,2}$ with compact support but not in $W^{1,p}$.  A node
carrying the printed class would therefore exclude the paper's own
substitutions and could not support the paper's use of (3.2) of the thesis.  For
$1<p<2$ the $W^{1,p}$ requirement would instead be redundant, since compact
support together with $L^2$ gives $L^p$ by H\"older on a set of finite measure;
so the clause binds exactly on the range where it cannot be discharged, and
Theorems~Theorem~1 of the thesis and~Theorem~2 of the thesis are needed for $p>2$.

Nothing is lost at the other end: Step 4 of the proof below shows that
(1.2) of the thesis itself extends from its printed $W^{1,p}$ test class to every
$\varphi\in W^{1,2}(\Omega)$ with compact support in $\Omega$, so the enlarged
class is available to the equation being differentiated as well.  The
exponent $2$ is exactly right for the display: the factor $D_jD_ku$ lies in
$L^2_{\mathrm{loc}}(\Omega)$, and the coefficient $D_{h_j}A^i(\nabla u)$ is
locally bounded because $\varepsilon>0$ and $\nabla u$ is continuous, so the
integrand is locally integrable as soon as $D_i\varphi\in L^2_{\mathrm{loc}}$.
\end{helper}
\begin{proof}
Write $e_k$ for the $k$-th standard basis vector of $\R^n$ and
$\Phi^i=A^i\circ g$ for the $i$-th component of the composite
$x\mapsto A(\nabla u(x))$.  Fix $k$, $\varphi$ and $g_\varphi$ as in the
statement and put $K_0=\operatorname{supp}\varphi$, a compact subset of
$\Omega$.

If $\Omega=\emptyset$ the asserted integral is an integral over the empty set
and vanishes, so assume $\Omega\ne\emptyset$ and fix once and for all a point
$x_\ast\in\Omega$.  By the differentiability clause of \noderef{EllipticStructureQ12} applied at
$x_\ast$ and $z=0$, the map $A$ is of class $C^1$ on
$D(\varepsilon)=\{h:\varepsilon+|h|^2\ne0\}$, which is all of $\R^n$ here
because $\varepsilon>0$; so $A$ is of class $C^1$ on $\R^n$.  This is the only
use of that clause and it is used in Step 2.  Conditions (Q1) of the thesis and
(Q2) of the thesis themselves are not used directly anywhere below: they enter only
as the structure hypothesis that \noderef{CitedRegularityEpsPos} requires in
order to yield Step 1, which is exactly how the cited result is stated.

\emph{Step 1: what (C1) supplies.}  Apply
\noderef{CitedRegularityEpsPos} with the present
$n,p,\varepsilon,\gamma_0,\gamma_1,\Omega,A$: its hypotheses are $\varepsilon>0$,
\noderef{EllipticStructureQ12} for the field $(x,z,h)\mapsto A(h)$ on $\Omega$,
and that $u$ with weak gradient $g$ is a weak solution of
$\operatorname{div}A(\nabla u)=0$ in $\Omega$ (\noderef{IsWeakSolutionOn}) --- all
three are hypotheses of the present lemma.  Its conclusion gives
\begin{enumerate}
\item $u\in W^{1,2}_{\mathrm{loc}}(\Omega)$ with weak gradient $g$
  (\noderef{MemW1pLocOn});
\item a field $g'$ with $g'$ continuous on $\Omega$ and $g=g'$ almost everywhere
  on $\Omega$.
\end{enumerate}
\emph{Item (2) does not assert continuity of the given $g$, and cannot: a weak
gradient is an almost-everywhere class, so continuity is a property of a
representative.  Every \emph{pointwise} consequence drawn below is therefore
drawn for $g'$, and transported back to $g$ across $g=g'$ a.e.\ wherever the
quantity it feeds is an integral or an almost-everywhere assertion, both of
which are blind to a null set.}

From (2) we get the local boundedness of the family of components \emph{of
$g'$}: if $K\subseteq\Omega$ is compact then $g'$ restricted to $K$ is
continuous, so $x\mapsto|g'(x)|$ attains a finite maximum $L$ on $K$ (a
continuous real function on a compact set is bounded), and
$|g'^j(x)|\le|g'(x)|\le L$ for every $x\in K$ and every $j$.  (If
$K=\emptyset$ any $L$ serves.)  The corresponding assertion for the given $g$ is
unavailable, and it is not needed: the only hypothesis in this proof that asks
for a bound at \emph{every} point of a compact set is the locally uniformly
bounded clause of \noderef{WeakGradientVectorComp} in Step 2, and Step 2 is fed
$g'$.

We also record, for use in Step 2, that each $g'^j$ lies in
$W^{1,2}_{\mathrm{loc}}(\Omega)$ with the \emph{same} weak gradient $Dg^j$ that
the lemma's hypothesis gives for $g^j$.  By \noderef{MemW1pLocOn} this is two
assertions, and each is an integral one.  First the identity of
\noderef{HasWeakGradientOn}: for every $\psi$ in that node's test class and every
index $i$, $\int_\Omega g^j\,D_i\psi\,dx=-\int_\Omega(Dg^j)_i\,\psi\,dx$;
since $g^j=g'^j$ almost everywhere on $\Omega$ the two left-hand integrands
agree almost everywhere, so the integrals are equal and
$\int_\Omega g'^j\,D_i\psi\,dx=-\int_\Omega(Dg^j)_i\,\psi\,dx$.  Second the
$L^2$ clauses: for compact $K\subseteq\Omega$, $g'^j=g^j$ almost everywhere on
$K$ and $Dg^j$ is unchanged, so $g'^j\in L^2(K)$ and $Dg^j\in L^2(K)$ follow
from the same statements for $g^j$.  Hence
\[
g'^j\in W^{1,2}_{\mathrm{loc}}(\Omega)\ \text{ with weak gradient }\ Dg^j
\qquad(j=1,\dots,n).
\tag{0}
\]

\emph{Step 2: the composite $A^i\circ\nabla u$ is in $W^{1,2}_{\mathrm{loc}}$.}
Fix $i$.  Identify $\R^m$ with $m=n$ and $\R^n$ in the usual way and let
$F=F^i\colon\R^n\to\R$ be the $i$-th component of $A$, i.e.\
$F(y)=A^i(y)$.  Then $F$ is of class $C^1$: $A$ is $C^1$ by the paragraph
above, the identification of $\R^n$ with the Euclidean space carrying $A$'s
argument is a linear homeomorphism and hence $C^\infty$, and the $i$-th
coordinate functional is linear and continuous, hence $C^\infty$; so
$F$ is a composition of $C^1$ and $C^\infty$ maps.

Now apply \noderef{WeakGradientVectorComp} with $m=n$, $q=2$, the open set
$\Omega$, the family $G_j=g'^j$, the fields $DG_j=Dg^j$, and this $F$.  Its
hypotheses hold: $1\le n$ and $1\le2$; $\Omega$ is open; for every $j$ we have
$g'^j\in W^{1,2}_{\mathrm{loc}}(\Omega)$ with weak gradient $Dg^j$
(\noderef{MemW1pLocOn}), which is $(0)$ of Step 1; the family $\{g'^j\}$ is
locally uniformly bounded by Step 1; and $F$ is $C^1$.  Its conclusion is
\[
\Phi'^i=F(g'^1,\dots,g'^n)\in W^{1,2}_{\mathrm{loc}}(\Omega)
\quad\text{with weak gradient}\quad
V'^i(x)=\sum_{j=1}^nD_jF\bigl(g'(x)\bigr)\,Dg^j(x) .
\]
Transport this back to $g$.  Both $\Phi^i=F(g^1,\dots,g^n)=A^i(g(\cdot))$ and
$V^i(x)=\sum_jD_jF(g(x))\,Dg^j(x)$ are obtained from $g$ by the same formulas
that produce $\Phi'^i$ and $V'^i$ from $g'$, so each pair agrees at every point
where $g=g'$, hence almost everywhere on $\Omega$ by Step 1.  Repeating the
unfolding of $(0)$ verbatim with $(\Phi'^i,V'^i)$ in place of $(g^j,Dg^j)$ ---
both sides of the \noderef{HasWeakGradientOn} identity and both $L^2$ clauses of
\noderef{MemW1pLocOn} are integrals, so each is unchanged when its integrand is
altered on a null set --- the membership passes to the unprimed pair:
\[
\Phi^i=F(g^1,\dots,g^n)\in W^{1,2}_{\mathrm{loc}}(\Omega)
\quad\text{with weak gradient}\quad
V^i(x)=\sum_{j=1}^nD_jF\bigl(g(x)\bigr)\,Dg^j(x) ,
\tag{1}
\]
and since $D_jF(y)$ is the $j$-th partial derivative of the $i$-th component of
$A$, that is $D_jF(g(x))=D_{h_j}A^i(\nabla u(x))$, the $k$-th component of
$V^i$ is
\[
(V^i)_k(x)=\sum_{j=1}^nD_{h_j}A^i(\nabla u(x))\,(Dg^j)_k(x) .
\tag{2}
\]
This is the step the $C^1$ clause buys: with bare differentiability of $A$ the
hypothesis ``$F$ is $C^1$'' of \noderef{WeakGradientVectorComp} is unavailable.

\emph{Step 3: a weak gradient vanishes off the support.}  We record the
following, for an arbitrary pair, because it is used twice below.  Let
$\psi_0\in W^{1,2}(\Omega)$ have weak gradient $g_0$ (\noderef{MemW1pOn} at the
exponent $2$) and compact support with
$L=\operatorname{supp}\psi_0\subseteq\Omega$.  Then
\[
g_0=0\qquad\text{almost everywhere on }\Omega\setminus L .
\tag{3}
\]
Indeed, $U=\Omega\setminus L$ is open, being the intersection of the open
$\Omega$ with the complement of the closed $L$.  Let $\psi\in C^\infty(\R^n)$
have compact support with $\operatorname{supp}\psi\subseteq U$, and fix an index
$i$.  Then $\psi$ is also admissible for $\Omega$, and the defining identity of
\noderef{HasWeakGradientOn} for the pair $(\psi_0,g_0)$ on $\Omega$, available
from \noderef{MemW1pOn}, gives
$\int_\Omega\psi_0\,D_i\psi=-\int_\Omega(g_0)_i\,\psi$.  The left-hand side is
$0$: the integrand vanishes off $\operatorname{supp}\psi$, and on
$\operatorname{supp}\psi\subseteq U$ we have $\psi_0=0$, because $\psi_0$
vanishes at every point outside $L=\operatorname{supp}\psi_0$.  Hence
$\int_{\R^n}\psi\,(g_0)_i=0$ for every such $\psi$, the integral over $\Omega$
and over $\R^n$ agreeing because the integrand vanishes off
$\operatorname{supp}\psi\subseteq U\subseteq\Omega$.  Moreover $(g_0)_i$ is
locally integrable on $U$: for compact $K'\subseteq U\subseteq\Omega$ we have
$|g_0|\in L^2(K')$ by \noderef{MemW1pOn} at exponent $2$, hence
$|g_0|\in L^1(K')$ by H\"older on a set of finite measure (Preamble), and
$|(g_0)_i|\le|g_0|$.  The fundamental lemma of the calculus of variations
(Preamble), whose test class is exactly the one just used, therefore gives
$(g_0)_i=0$ almost everywhere on $U$.  As $i$ was arbitrary, (3) follows.

\emph{Step 4: (1.2) of the thesis extends to the $W^{1,2}$ test class with compact
support.}  We claim that for every pair $(\psi_0,g_0)$ as in Step 3, that is
every $\psi_0\in W^{1,2}(\Omega)$ with weak gradient $g_0$
(\noderef{MemW1pOn}) and compact support
$L=\operatorname{supp}\psi_0\subseteq\Omega$,
\[
\int_\Omega\sum_{i=1}^n\Phi^i\,(g_0)_i\,dx=0 .
\tag{4}
\]
This is the only place where the enlarged test class of the present statement is
earned, and it is the reason the $W^{1,p}$ hypothesis is not needed.

If $L=\emptyset$ then $\psi_0$ vanishes identically, so $g_0=0$ almost
everywhere on $\Omega$ by (3) and (4) is immediate.  Assume $L\ne\emptyset$.
Apply \noderef{WeakGradientSmoothApprox} with this $n$, the exponent $q=2$, the
open set $\Omega$ and the pair $(\psi_0,g_0)$: its hypotheses are $1\le n$
(from $n\ge2$), $1\le2$, openness of $\Omega$, $\psi_0\in W^{1,2}(\Omega)$ with
weak gradient $g_0$, compactness of the support and
$\operatorname{supp}\psi_0\subseteq\Omega$, all of which hold.  It produces a
compact $K\subseteq\Omega$ with $L\subseteq K$ and functions
$\Psi_m\in C^\infty(\R^n)$, each with compact support contained in $K$, such
that
\[
\int_K|\Psi_m-\psi_0|^2dx\longrightarrow0
\qquad\text{and}\qquad
\int_K\bigl|D_i\Psi_m-(g_0)_i\bigr|^2dx\longrightarrow0
\quad(m\to\infty)
\tag{5}
\]
for every index $i$.

Each $\Psi_m$ is admissible in (1.2) of the thesis.  Indeed
\noderef{SmoothCompactSupportMemW1p} applies with the exponent $r=p$ --- its
hypotheses are $1\le n$, $1\le p$ (from $p>1$), openness of $\Omega$,
smoothness of $\Psi_m$, compactness of its support and
$\operatorname{supp}\Psi_m\subseteq K\subseteq\Omega$ --- and gives
$\Psi_m\in W^{1,p}(\Omega)$ with weak gradient the classical gradient
$\nabla\Psi_m$ (\noderef{MemW1pOn}).  So the weak formulation of
\noderef{IsWeakSolutionOn}, whose test class is ``$W^{1,p}(\Omega)$ with weak
gradient, compact support, support inside $\Omega$'' and whose inhomogeneity is
$B\equiv0$, gives
\[
\int_\Omega\sum_{i=1}^n\Phi^i\,(D_i\Psi_m)\,dx=0\qquad\text{for every }m .
\tag{6}
\]

Now pass to the limit.  First, $\Phi$ is bounded almost everywhere on $K$, and
that is all the estimate below uses.  By Step 1 the map $g'$ is continuous on
$\Omega$, so $|g'|$ attains a finite maximum $R_K$ on the compact
$K\subseteq\Omega$; the closed ball $\{y:|y|\le R_K\}$ is compact and $A$ is
continuous, being $C^1$, so $|A|$ is bounded on it, and therefore
$\Phi'^i=A^i(g'(\cdot))$ satisfies
$\sup_{x\in K}\max_i|\Phi'^i(x)|<\infty$.  Since $g=g'$ almost everywhere on
$\Omega$ we have $\Phi^i=\Phi'^i$ almost everywhere on $\Omega$, so there is a
finite $M$ with
\[
\max_{1\le i\le n}|\Phi^i(x)|\le M\qquad\text{for almost every }x\in K ,
\]
which dominates exactly as an everywhere bound would in the integral estimate
below.  Each $\Phi^i$ is also measurable on $K$: $\Phi'^i$ is continuous there,
being the composition of the continuous $A^i$ with the continuous $g'$, and
$\Phi^i=\Phi'^i$ almost everywhere, so $\Phi^i$ is almost everywhere equal on
$K$ to a continuous function.  Second, both integrals in (6) and (4) are
integrals over $K$: each $D_i\Psi_m$ vanishes off $\operatorname{supp}\Psi_m
\subseteq K$, because $\Psi_m$ vanishes identically on the open complement of
its support and a function constant on an open set has vanishing derivative
there; and $(g_0)_i=0$ almost everywhere on $\Omega\setminus K$ by (3),
since $\Omega\setminus K\subseteq\Omega\setminus L$.  Third, all the integrals
below are finite: $\Phi^i$ is bounded on $K$, each $D_i\Psi_m$ is continuous
hence bounded, and $(g_0)_i\in L^2(K)\subseteq L^1(K)$ by H\"older on a set of
finite measure (Preamble), $|K|<\infty$.  Therefore
\[
\Bigl|\int_K\sum_{i=1}^n\Phi^i\bigl[D_i\Psi_m-(g_0)_i\bigr]dx\Bigr|
\le M\sum_{i=1}^n\int_K\bigl|D_i\Psi_m-(g_0)_i\bigr|dx
\le M\sum_{i=1}^n|K|^{1/2}\Bigl(\int_K\bigl|D_i\Psi_m-(g_0)_i\bigr|^2dx\Bigr)^{1/2},
\]
the second inequality being H\"older on the finite-measure set $K$ with the
conjugate exponents $2$ and $2$ (Preamble).  By (5) each of the $n$ summands
tends to $0$, so the left-hand side does.  Combining with (6), whose left-hand
side is $0$ for every $m$ and equals $\int_K\sum_i\Phi^iD_i\Psi_m$, we get
$\int_K\sum_i\Phi^i(g_0)_i=0$, hence (4).  This proves the claim.

\emph{Step 5: geometry and the cut-off.}  If $K_0=\emptyset$ then $\varphi$
vanishes identically, so $g_\varphi=0$ almost everywhere on $\Omega$ by (3),
the asserted integrand vanishes almost everywhere, and the integral is $0$.
Assume therefore $K_0\ne\emptyset$.  The set $K_0$ is compact,
$\R^n\setminus\Omega$ is closed, and the two are disjoint, so
$d=\operatorname{dist}(K_0,\R^n\setminus\Omega)>0$ (with $d=+\infty$ when
$\Omega=\R^n$).  Fix $\delta>0$ with $4\delta<d$ and for $s>0$ put
$K_s=\{x:\operatorname{dist}(x,K_0)\le s\}$, which is closed and bounded, hence
compact, and satisfies $K_s\subseteq\Omega$ for $s\le3\delta$.  By (3) applied
to $(\varphi,g_\varphi)$, whose support is $K_0$, we have $g_\varphi=0$ almost
everywhere on $\Omega\setminus K_0$; put
\[
\widehat g_\varphi=g_\varphi\,\chi_{K_0},
\]
so that $\widehat g_\varphi=g_\varphi$ almost everywhere on $\Omega$ and
$\widehat g_\varphi$ vanishes everywhere off $K_0$.

\emph{The cut-off.}  Fix
$\rho\in C^\infty(\R^n)$ with $\rho\ge0$,
$\operatorname{supp}\rho\subseteq\overline{B_1(0)}$ and $\int_{\R^n}\rho=1$ (the
standard choice $\rho(z)=c\exp(1/(|z|^2-1))$ for $|z|<1$ and $0$ otherwise), and
put $\rho_\delta(z)=\delta^{-n}\rho(z/\delta)$, so that $\rho_\delta\in
C^\infty$, $\rho_\delta\ge0$,
$\operatorname{supp}\rho_\delta\subseteq\overline{B_\delta(0)}$ and
$\int\rho_\delta=1$.  Define
\[
\chi=\chi_{K_{2\delta}}*\rho_\delta,
\qquad
\chi(x)=\int_{\R^n}\chi_{K_{2\delta}}(y)\rho_\delta(x-y)\,dy .
\]
Then: $\chi\in C^\infty(\R^n)$, by the Preamble item on convolution with a
smooth compactly supported kernel, $\chi_{K_{2\delta}}$ being bounded with
compact support and hence locally integrable.  $0\le\chi\le1$, because
$\rho_\delta\ge0$ integrates to $1$ and $0\le\chi_{K_{2\delta}}\le1$.
$\chi=1$ on the open set $\{x:\operatorname{dist}(x,K_0)<\delta\}$: for such
$x$ and $|x-y|\le\delta$ one has
$\operatorname{dist}(y,K_0)\le|x-y|+\operatorname{dist}(x,K_0)<2\delta$, so
$y\in K_{2\delta}$ and the integrand is $\rho_\delta(x-y)$, whose integral is
$1$.  Consequently $\chi$ is constant on that open set and all its partial
derivatives vanish there.  Finally $\chi=0$ off $K_{3\delta}$: if
$\operatorname{dist}(x,K_0)>3\delta$ and $|x-y|\le\delta$ then
$\operatorname{dist}(y,K_0)>2\delta$, so $y\notin K_{2\delta}$ and the
integrand vanishes.  Hence $\operatorname{supp}\chi\subseteq
K_{3\delta}\subseteq\Omega$ is compact.

\emph{Step 6: the translated test function and the exact difference identity.}
Fix $h$ with $0<|h|\le\delta$ and set
\[
\varphi_h=\varphi(\cdot-he_k),
\qquad
G_h=\widehat g_\varphi(\cdot-he_k).
\]
We claim $(\varphi_h,G_h)$ satisfies the hypotheses of Step 4, i.e.\ that
$\varphi_h\in W^{1,2}(\Omega)$ with weak gradient $G_h$
(\noderef{MemW1pOn}), $\varphi_h$ has compact support and
$\operatorname{supp}\varphi_h\subseteq\Omega$.

\emph{Support.}  $\varphi_h$ vanishes off $K_0+he_k$, which is compact, and
$\operatorname{dist}(x+he_k,K_0)\le|h|\le\delta$ gives
$K_0+he_k\subseteq K_\delta\subseteq\Omega$.  Likewise $G_h$ vanishes off
$K_0+he_k$.

\emph{The $L^2$ memberships.}  The functions $\varphi$ and $g_\varphi$ are
almost everywhere strongly measurable on $\Omega$ by \noderef{MemW1pOn}, and
translation preserves Lebesgue measurability (Preamble), so $\varphi_h$ and
$G_h$ are measurable.  By translation invariance of Lebesgue measure
(Preamble) and because the integrands vanish off $K_0+he_k\subseteq\Omega$ resp.\
$K_0\subseteq\Omega$,
\[
\int_\Omega|\varphi_h|^2\,dx=\int_{\R^n}|\varphi(x-he_k)|^2\,dx
=\int_{\R^n}|\varphi|^2\,dy=\int_\Omega|\varphi|^2\,dy<\infty,
\]
the last integral finite by \noderef{MemW1pOn} at exponent $2$; and the same
computation with $|\widehat g_\varphi|\le|g_\varphi|$ in place of $\varphi$ gives
$\int_\Omega|G_h|^2<\infty$.

\emph{The weak-gradient identity.}  Let $\psi\in C^\infty(\R^n)$ have compact
support with $\operatorname{supp}\psi\subseteq\Omega$ and fix an index $i$.  Put
$\widehat\psi=\psi(\cdot+he_k)$, again $C^\infty$ with compact support, and note
$D_i\widehat\psi=(D_i\psi)(\cdot+he_k)$.  Then
\[
\int_\Omega\varphi_h\,D_i\psi\,dx
=\int_{\R^n}\varphi(x-he_k)(D_i\psi)(x)\,dx
=\int_{\R^n}\varphi(y)\,(D_i\psi)(y+he_k)\,dy
=\int_{\R^n}\varphi\,D_i\widehat\psi\,dy ,
\]
the first equality because the integrand vanishes off $K_0+he_k\subseteq\Omega$
and the second by translation invariance (Preamble).  Next,
$\varphi$ vanishes off $K_0$, and on the open set
$\{\operatorname{dist}(\cdot,K_0)<\delta\}\supseteq K_0$ we have $\chi=1$ and
$D_i\chi=0$ by Step 5, so there
$D_i(\chi\widehat\psi)=\chi D_i\widehat\psi+\widehat\psi D_i\chi
=D_i\widehat\psi$; hence the two integrands $\varphi D_i\widehat\psi$ and
$\varphi D_i(\chi\widehat\psi)$ agree everywhere and
\[
\int_{\R^n}\varphi\,D_i\widehat\psi
=\int_{\R^n}\varphi\,D_i\bigl(\chi\widehat\psi\bigr)
=\int_\Omega\varphi\,D_i\bigl(\chi\widehat\psi\bigr) ,
\]
the last equality because the integrand vanishes off $K_0\subseteq\Omega$.  The
function $\chi\widehat\psi$ is $C^\infty$, has compact support contained in
$\operatorname{supp}\chi$, and $\operatorname{supp}(\chi\widehat\psi)\subseteq
\operatorname{supp}\chi\subseteq\Omega$, so it is admissible in
\noderef{HasWeakGradientOn} for $\Omega$.  Applying the defining identity for
$(\varphi,g_\varphi)$, available from \noderef{MemW1pOn},
\[
\int_\Omega\varphi\,D_i\bigl(\chi\widehat\psi\bigr)
=-\int_\Omega(g_\varphi)_i\,\chi\widehat\psi
=-\int_\Omega(\widehat g_\varphi)_i\,\widehat\psi ,
\]
the last equality because $g_\varphi=\widehat g_\varphi$ almost everywhere on
$\Omega$ by Step 5, $\widehat g_\varphi$ vanishes off $K_0$, and $\chi=1$ on
$K_0$.  Undoing the translation exactly as above,
\[
-\int_\Omega(\widehat g_\varphi)_i\,\widehat\psi
=-\int_{\R^n}(\widehat g_\varphi)_i(y)\,\psi(y+he_k)\,dy
=-\int_{\R^n}(G_h)_i(x)\,\psi(x)\,dx
=-\int_\Omega(G_h)_i\,\psi ,
\]
using translation invariance again and that the integrands vanish off
$K_0\subseteq\Omega$ resp.\ $K_0+he_k\subseteq\Omega$.  Chaining the four
displays gives
$\int_\Omega\varphi_h D_i\psi=-\int_\Omega(G_h)_i\psi$ for every admissible
$\psi$ and every $i$, which is \noderef{HasWeakGradientOn} for
$(\varphi_h,G_h)$ on $\Omega$.  Together with the $L^2$ memberships this proves
the claim.

Now apply (4) of Step 4 twice: to the pair $(\varphi,g_\varphi)$, which lies in
$W^{1,2}(\Omega)$ with weak gradient $g_\varphi$ and compact support
$K_0\subseteq\Omega$ by hypothesis, and to the pair $(\varphi_h,G_h)$ just
shown to satisfy the same hypotheses with support $K_0+he_k\subseteq\Omega$.
This gives
\[
\int_\Omega\sum_{i=1}^n\Phi^i(x)\,(g_\varphi)_i(x)\,dx=0,
\qquad
\int_\Omega\sum_{i=1}^n\Phi^i(x)\,(G_h)_i(x)\,dx=0 .
\]
In the first integral replace $g_\varphi$ by $\widehat g_\varphi$ (they agree
almost everywhere on $\Omega$ by Step 5) and restrict to $K_0$, off which
$\widehat g_\varphi$ vanishes; in the second substitute $x=y+he_k$ by
translation invariance (Preamble) and restrict to $K_0$, off which $G_h(y+he_k)
=\widehat g_\varphi(y)$ vanishes.  We obtain
\[
\int_{K_0}\sum_i\Phi^i(y)\,(\widehat g_\varphi)_i(y)\,dy=0,
\qquad
\int_{K_0}\sum_i\Phi^i(y+he_k)\,(\widehat g_\varphi)_i(y)\,dy=0 .
\]
Both integrals are finite: $K_\delta\subseteq\Omega$ is compact, so $g'$ is
bounded on $K_\delta$ by Step 1 and $A$ is continuous, whence each $\Phi'^i$ is
bounded on $K_\delta$, and $\Phi^i=\Phi'^i$ almost everywhere gives the same
bound for $\Phi^i$ almost everywhere on $K_\delta$, which is what the
integrals need; and $(\widehat g_\varphi)_i\in L^1(K_0)$ because
$|g_\varphi|\in L^2(K_0)\subseteq L^1(K_0)$ by H\"older on a set of finite
measure (Preamble).  For
$y\in K_0$ we have $y+he_k\in K_\delta$, so the second integrand is bounded as
well.  Subtracting and dividing by $h\ne0$,
\[
\int_{K_0}\sum_{i=1}^n\frac{\Phi^i(y+he_k)-\Phi^i(y)}{h}\,(\widehat
g_\varphi)_i(y)\,dy=0
\qquad(0<|h|\le\delta).
\tag{7}
\]

\emph{Step 7: passing to the limit $h\to0$.}  Fix $i$.  By Step 2 we have
$\Phi^i\in W^{1,2}_{\mathrm{loc}}(\Omega)$ with weak gradient $V^i$, and
$K_0\subseteq\Omega$ is compact, so
\noderef{WeakGradientDifferenceQuotient} applies with $q=2$, the open set
$\Omega$, the function $\Phi^i$, its weak gradient $V^i$, the compact set $K_0$
and the index $k$, giving
\[
\int_{K_0}\Bigl|\frac{\Phi^i(y+he_k)-\Phi^i(y)}{h}-(V^i)_k(y)\Bigr|^2dy
\longrightarrow0\qquad(h\to0,\ h\ne0).
\]
Also $(\widehat g_\varphi)_i\in L^2(K_0)$: indeed $|g_\varphi|\in L^2(\Omega)$ by
\noderef{MemW1pOn} at exponent $2$, and $|(\widehat g_\varphi)_i|\le|g_\varphi|$.
Hence by the Cauchy--Schwarz inequality for the integral, which is H\"older with
the conjugate exponents $2$ and $2$ on $K_0$ (Preamble),
\[
\Bigl|\int_{K_0}\sum_i\Bigl[\frac{\Phi^i(\cdot+he_k)-\Phi^i}{h}-(V^i)_k\Bigr]
(\widehat g_\varphi)_i\Bigr|
\le\sum_i\Bigl\|\frac{\Phi^i(\cdot+he_k)-\Phi^i}{h}-(V^i)_k\Bigr\|_{L^2(K_0)}
\bigl\|(\widehat g_\varphi)_i\bigr\|_{L^2(K_0)},
\]
a finite sum of products each of which tends to $0$.  Combining with (7), whose
left-hand side is $0$ for every $h$ with $0<|h|\le\delta$, and letting $h\to0$
along such $h$,
\[
\int_{K_0}\sum_{i=1}^n(V^i)_k\,(\widehat g_\varphi)_i\,dy=0 .
\]
The integrand vanishes off $K_0$ and $\widehat g_\varphi=g_\varphi$ almost
everywhere on $\Omega$, so this says
\[
\int_\Omega\sum_{i=1}^n(V^i)_k\,(g_\varphi)_i\,dx=0 .
\tag{8}
\]

\emph{Step 8: restoring the thesis' index order.}  By Step 1 we have
$u\in W^{1,2}_{\mathrm{loc}}(\Omega)$ with weak gradient $g$, and by hypothesis
each $g^i\in W^{1,2}_{\mathrm{loc}}(\Omega)$ with weak gradient $Dg^i$
(\noderef{MemW1pLocOn}); $\Omega$ is open, $1\le n$ and $1\le2$.  So
\noderef{WeakSecondDerivativeSymmetry} applies at $q=2$ and gives, for each pair
$(j,k)$,
\[
(Dg^k)_j=(Dg^j)_k\qquad\text{almost everywhere on }\Omega .
\]
There are finitely many pairs $(j,k)$, so the union of the corresponding null
sets is null and all $n^2$ identities hold simultaneously almost everywhere on
$\Omega$.  Substituting into (2), for almost every $x\in\Omega$
\[
\sum_{i=1}^n(V^i)_k(x)\,(g_\varphi)_i(x)
=\sum_{i=1}^n\sum_{j=1}^nD_{h_j}A^i(\nabla u(x))\,(Dg^k)_j(x)\,(g_\varphi)_i(x) ,
\]
and two integrands agreeing almost everywhere have the same integral.  With (8)
this is exactly the assertion
\[
\int_\Omega\Bigl\{\sum_{i,j=1}^nD_{h_j}A^i(\nabla u)\,D_jD_ku\,D_i\varphi\Bigr\}dx=0 ,
\]
since $D_jD_ku=(Dg^k)_j$ and $D_i\varphi=(g_\varphi)_i$.
\end{proof}
